\documentclass{IOS-Book-Article}

\usepackage[utf8]{inputenc}
\usepackage[T1]{fontenc}
\usepackage{mathptmx}
\usepackage{graphicx}
\usepackage{booktabs}
\usepackage{hyperref}
\usepackage{url}
\raggedbottom

\def\hb{\hbox to 11.5 cm{}}

\begin{document}

\pagestyle{headings}
\def\thepage{}
\begin{frontmatter}

\title{TrashUQ: Federated Learning for Intelligent Waste Monitoring on Arduino UNO Q}

\markright{}{}

\author[a]{\fnms{Aleix} \snm{Bertran Andreu}}
\author[A]{\fnms{Pol} \snm{Llin\`as Vaquer}}
\author[A]{\fnms{Bru} \snm{Pall\`as Vargu\'es}}
\author[A]{\fnms{Aleix} \snm{Rosinach Olivart}}

\runningauthor{TrashUQ Team}

\address[A]{Universitat de Lleida, Lleida, Catalonia}

\begin{abstract}
This paper presents TrashUQ, a waste-monitoring platform that integrates Arduino UNO Q edge nodes, MQTT telemetry, and a gRPC Federated Learning (FL) coordinator. We report two validations. First, a real two-device deployment confirms end-to-end telemetry and FL round-trip behavior (282 persisted MQTT records, model version 0$\rightarrow$26, 25/30 successful aggregations). Second, a controlled non-IID FL simulation (Dirichlet $\alpha=0.3$) with 2, 5, 10, and 20 clients shows stable FedAvg convergence over 25 rounds, with final accuracies around 93--94\%. We further demonstrate that stochastic rounding gradient quantization achieves 4--16$\times$ model compression with accuracy loss below 2.84 p.p.\ at 20 clients, outperforming Top-k sparsification which either saves no communication or degrades accuracy by 9+~p.p. FL convergence is robust across a wide range of non-IID intensities ($\alpha = 0.1$--$1.0$). Results indicate that FL orchestration is feasible on this embedded stack, and that stochastic rounding is the preferred method for communication-constrained edge deployments.
\end{abstract}

\begin{keyword}
Federated Learning\sep Edge AI\sep Waste Classification\sep
Arduino UNO Q\sep MQTT\sep TrashNet
\end{keyword}
\end{frontmatter}
\markboth{TrashUQ: FL for Waste Monitoring on Arduino UNO Q}{TrashUQ: FL for Waste Monitoring on Arduino UNO Q}

\section{Introduction}
\label{sec:introduction}
% This file is input by the main document. Do not add \begin or \end.

Waste mismanagement poses critical environmental and public health
challenges. Intelligent waste sorting systems powered by deep learning can
significantly improve recycling rates, but conventional approaches rely on
centralized cloud infrastructure that raises privacy, latency, and bandwidth
concerns~\cite{mcmahan2017fedavg}. Federated Learning (FL) addresses these
limitations by training models collaboratively across edge devices while
keeping data local~\cite{li2020fedavg,kairouz2021fl}.

Several compression strategies have been proposed to mitigate the
communication bottleneck in FL, including stochastic quantization
(QSGD~\cite{alistarh2017qsgd}), 1-bit sign compression
(SignSGD~\cite{bernstein2018signsgd}), top-k sparsification~\cite{aji2017topk},
and unified compressed FL frameworks (FedCOM~\cite{haddadpour2021fedcom}).
Hierarchical FL~\cite{liu2020hierarchical,abad2020hierarchical} further
reduces coordinator load through multi-level aggregation. However, most prior
work evaluates these techniques exclusively in simulation with high-end
hardware, leaving open questions about their applicability to resource-
constrained edge deployments. The Arduino UNO Q---featuring a Qualcomm
QRB2210 processor---enables practical FL deployment on physical low-cost
devices.

This paper presents TrashUQ, a complete FL-based waste monitoring system
deployed on two Arduino UNO Q boards. Our contributions are:
\begin{enumerate}
  \item A real two-node FL deployment with MQTT telemetry validation,
  confirming end-to-end FL round-trip behavior (282 records, 25/30 successful
  aggregations).
  \item FL simulation results on a controlled synthetic dataset with FedAvg
  across 2--20 non-IID clients, demonstrating stable convergence around
  93--94\% accuracy.
  \item Stochastic rounding gradient quantization achieving 4--8$\times$ model
  compression with negligible accuracy degradation (0--1.33\%).
  \item Open-source implementation at
  \url{https://github.com/TrashUQ/TrashUQ}~\cite{trashuqrepo}.
\end{enumerate}


\section{System Architecture}
\label{sec:architecture}
\subsection{Hardware Platform}

Each Arduino UNO Q has a Qualcomm QRB2210 application processor (Debian Linux) and an STM32U585 microcontroller. Boards connect to the backend over Wi-Fi.

\subsection{Edge Software Pipeline}

The edge daemon follows capture/classify/act logic. Classification uses a MobileNetV2 TFLite model ($224\times224$, classes: cardboard, glass, paper, plastic). User corrections are stored locally and used for on-device fine-tuning.

\subsection{Federated Learning Pipeline}

The trainable model is a frozen MobileNetV2 backbone with a small learnable head. After accumulating new labels, a node starts an FL round: it fetches global weights, trains locally, and submits an update via gRPC. The coordinator applies FedAvg~\cite{mcmahan2017fedavg}.

\subsection{MQTT Telemetry and Backend}

Boards publish status/events/FL metrics to \texttt{arduino/<device-id>/\#}. A FastAPI backend persists telemetry in PostgreSQL and exposes the FL coordinator. A Next.js dashboard provides real-time monitoring via MQTT over WebSocket.
Code and deployment artifacts are available at
\url{https://github.com/TrashUQ/TrashUQ}.


\section{Hardware Validation}
\label{sec:hardware}
The hardware validation was executed on May 16, 2026, using two Arduino UNO Q nodes connected to a shared backend stack (MQTT broker, PostgreSQL, REST API, dashboard, and gRPC FL coordinator). Both nodes ran on Debian 13 and used the same edge daemon, topic contract, and FL client API.

\subsection{Protocol Validation and Persistence}
We validated end-to-end ingestion for both devices across the path Edge $\rightarrow$ MQTT $\rightarrow$ Backend $\rightarrow$ PostgreSQL $\rightarrow$ Frontend. During a $\sim$12-minute run, the system persisted 282 MQTT rows in total: 147 from \texttt{unoq-01} and 135 from \texttt{unoq-02}. Per-device topic counts were 27/21 status messages, 16/14 metric messages, and 104/100 event messages, respectively.

\subsection{FL Round-Trip Behavior on Real Nodes}
Both nodes successfully completed FL handshake and training-update cycles with the coordinator (\texttt{Join}, \texttt{GetGlobalModel}, \texttt{SubmitUpdate}). The experiment reached model version 26, with 25 successful aggregations out of 30 submissions (83.3\% aggregation success). Five submissions were rejected as stale during concurrent updates, confirming that coordinator-side monotonic round checks were enforced correctly.

Round submission latency (from local training completion to update submission) was low (median $\sim$25\,ms, max 105\,ms). Mean local metrics were 0.921 accuracy / 0.236 loss for \texttt{unoq-01}, and 0.871 accuracy / 0.311 loss for \texttt{unoq-02}.

\subsection{Operational Limitations}
Both devices used \texttt{--fake-camera} and HTTP label injection to exercise FL deterministically. Therefore, camera-in-the-loop inference latency and classification-topic throughput were not measured in this run.


\section{FL Simulation and Communication Efficiency}
\label{sec:fl_simulation}
We ran FedAvg for 25 rounds with non-IID partitions (Dirichlet $\alpha=0.3$), using 3 seeds per setting with client counts \{2, 5, 10, 20\}. Each run used local epochs $=2$, batch size $=16$, SGD with learning rate $0.18$, and full participation.

To enable reproducible experiments without external dependencies, we generated a synthetic dataset of $16\times16$ grayscale images (4 waste classes, 220 samples/class). Each class prototype---horizontal stripes (cardboard), diagonal lines (glass), nested rectangles (paper), vertical stripes (plastic)---is blended via convex combination (52\% main, 23\% next, 10\% alternate, 15\% shared) with Gaussian noise ($\sigma=0.30$), brightness jitter, and pixel roll. The simulator can optionally load real TrashNet images if available.

\subsection{Convergence and Communication Efficiency}
Across all settings, FL converged stably with final accuracy near 93--94\% and no failed rounds (Table~\ref{tab:fl-summary}). Mean round duration increased from 0.005\,s (2 clients) to 0.008\,s (20 clients), while total communication grew approximately linearly from 0.394\,MB to 3.941\,MB.

\begin{table}[t]
\centering
\caption{FL simulation summary (mean across 3 seeds).}
\label{tab:fl-summary}
\begin{tabular}{lcccc}
\toprule
Clients & Acc. (\%) & Loss & Round Time (s) & Comm. (MB) \\
\midrule
2 & 93.75 & 0.2192 & 0.005 & 0.394 \\
5 & 93.56 & 0.2004 & 0.005 & 0.985 \\
10 & 93.37 & 0.2018 & 0.006 & 1.970 \\
20 & 93.75 & 0.1933 & 0.008 & 3.941 \\
\bottomrule
\end{tabular}
\end{table}

\subsection{Gradient Quantization}
\label{sec:quantization}

To address the communication bottleneck, we integrate stochastic rounding quantization into the FedAvg pipeline. After local training, each client quantizes weight updates before transmission; the server dequantizes before weighted averaging. We evaluate 8-bit (4$\times$ weight compression) and 4-bit with uint8 packing (8$\times$ weight compression).

\begin{table}[t]
\centering
\caption{Quantization results (20 clients, mean across 3 seeds).}
\label{tab:quantization}
\begin{tabular}{lcccc}
\toprule
Config & Acc. (\%) & Loss & Comm. (MB) & Compression \\
\midrule
float32 & 93.75 & 0.1933 & 3.941 & 1.0$\times$ \\
8-bit   & 93.75 & 0.1935 & 2.474 & $\sim$4$\times$ \\
4-bit   & 93.75 & 0.1976 & 2.229 & $\sim$8$\times$ \\
\bottomrule
\end{tabular}
\end{table}

The 8-bit configuration preserves accuracy identically (93.75\%) with $\sim$4$\times$ model compression. The 4-bit packed variant achieves $\sim$8$\times$ compression with a maximum accuracy degradation of 1.33\% (at 2 clients), diminishing to zero as the client count increases. These results demonstrate that lightweight gradient quantization effectively mitigates the communication bottleneck without compromising model quality on resource-constrained edge hardware.


\section{Conclusion and Future Work}
\label{sec:conclusion}
TrashUQ validates FL on real Arduino UNO Q hardware (two-node deployment, 282 MQTT records, 25/30 successful FL aggregations) and in controlled simulation (FedAvg, 2--20 non-IID clients, 93--94\% accuracy). Stochastic rounding achieves 4--16$\times$ compression with accuracy loss below 2.84~p.p.\ at 20 clients, outperforming Top-k sparsification which either saves no bytes or degrades accuracy by 9+~p.p. FL convergence is robust across non-IID intensities ($\alpha=0.1$ to $1.0$). By contrasting our flat topology with hierarchical FL and comparing stochastic rounding to sparsification, we position TrashUQ as a practical alternative requiring no hyperparameter tuning at the edge. Code at \url{https://github.com/TrashUQ/TrashUQ}; future work includes camera-in-the-loop validation and on-device quantized deployment.


\vspace*{-2mm}\section*{Acknowledgements}
The authors thank Jordi Planes Cid from the Universitat de Lleida for guidance and support during this work.

\vspace*{-4mm}
\bibliographystyle{vancouver}
\bibliography{references}

\end{document}
